\loai{Tính $U_2$}
\begin{vidu} %[3P3-14.2K][Loai1] 
	Một máy biến thế dùng làm máy giảm thế (hạ thế) gồm cuộn dây 100 vòng và cuộn dây 500 vòng. Bỏ qua mọi hao phí của máy biến thế. Khi nối hai đầu cuộn sơ cấp với hiệu điện thế $u = 100\sqrt{2}\sin100\pi  t$  V thì hiệu điện thế hiệu dụng ở hai đầu cuộn thứ cấp bằng
	\choice
	{10 V}
	{\True 20 V}
	{50 V}
	{500 V}
	\loigiai{
%\Binhluan{
		Máy giảm thế nên ${N_1}>{N_2}$ ta có:
		\begin{align*}
			& {N_1}=500;\ {N_2}=100;\  {U_1}=\dfrac{100\sqrt{2}}{\sqrt{2}}=100\ V \\
			&\dfrac{U_1}{U_2}=\dfrac{N_1}{N_2}\Leftrightarrow {U_2}=\dfrac{N_2}{N_1}. {U_1}=20\ V.
		\end{align*}
%}
%{Học sinh nhớ đổi sang giá trị hiệu dụng rồi áp dụng công thức}
	}
\end{vidu}
\begin{vidu} %[3P3-14.2B][Loai1] 
	Một máy biến thế có số vòng dây của cuộn sơ cấp là 5000 vòng và thứ cấp là 1000 vòng. Bỏ qua mọi hao phí qua máy biến thế. Đặt vào hai đầu cuộn sơ cấp hiệu điện thế xoay chiều có giá trị hiệu dụng 100 V thì hiệu điện thế ở hai đầu cuộn thứ cấp có giá trị cực đại là
	\choice
	{\True 28,3 V}
	{40 V}
	{20 V}
	{56,4 V}
	\loigiai{
%\Binhluan{
			Ta có: 
			\begin{align*}
			\dfrac{U_1}{U_2}=\dfrac{N_1}{N_2}&\Rightarrow U_2=U_1\dfrac{N_2}{N_1}=100.\dfrac{1000}{5000}=20\ V\\
			&\Rightarrow U_{02}=20\sqrt{2}\approx 28,28\ V.
			\end{align*}	
%}
%{Đề bài hỏi giá trị cực đại nên nếu "nhanh ẩu đoảng" sẽ bị nhầm sang đáp án C}
	}
\end{vidu}
\loai{Tính số vòng dây}
\begin{vidu} %[3P3-14.2B][Loai2] 
	Một máy biến thế có cuộn sơ cấp 1000 vòng dây được mắc vào mạng điện xoay chiều có hiệu điện thế hiệu dụng 220 V. Khi đó hiệu điện thế hiệu dụng ở hai đầu cuộn thứ cấp để hở là 484 V. Bỏ qua mọi hao phí của máy biến thế. Số vòng dây của cuộn thứ cấp là
	\choice
	{2500}
	{1100}
	{2000}
	{\True 2200}
	\loigiai{
	%	\Binhluan{
		\begin{align*}
			&{N_1}=1000;\ {U_1}=220\ V;\ {U_2}=484\ V\\ 
			&\dfrac{U_1}{U_2}=\dfrac{N_1}{N_2}\Rightarrow {N_2}=\dfrac{U_2}{U_1}{N_1}=2200\ \text{vòng}.	
		\end{align*}
	%}
%{
%Trong các kí hiệu máy biến áp thì số 1 ứng với sơ cấp. Số 2 ứng với thứ cấp.
%}
	}
\end{vidu}
\loai{Tính tỉ số số vòng dây}
\begin{vidu} %[3P3-14.2B][Loai3] 
	Một máy biến áp có điện trở các cuộn dây không đáng kể. Điện áp hiệu dụng giữa hai đầu cuộn thứ cấp và điện áp hiệu dụng giữa hai đầu cuộn sơ cấp lần lượt là 55 V và 220 V. Bỏ qua các hao phí trong máy, tỉ số giữa số vòng dây cuộn sơ cấp và số vòng dây cuộn thứ cấp bằng
	\choice
	{8}
	{\True 4}
	{2}
	{$\dfrac{1}{4}$}
	\loigiai{
%\Binhluan{
		\begin{align*}
			&{U_2}=55\ V ;\
			{U_1}=220\ V\\
			&\dfrac{U_1}{U_2}=\dfrac{N_1}{N_2}=4.	
		\end{align*}
%}
%{
%Tỉ số giữa số vòng dây 2 cuộn thứ cấp và sơ cấp cũng chính là tỉ số giữa các điện áp tương ứng.
%}
	}
\end{vidu}

\loai{Tính số vòng dây - cho biểu thức $\Phi$}
\begin{vidu} %[3P3-14.2B][Loai5] 
	Từ thông gửi qua một vòng dây của lõi sắt nằm trong cuộn sơ cấp một máy biến áp có dạng $\Phi_0=0,9\cos 100\pi t\ mWb$. Biết lõi sắt khép kín các đường sức từ. Nếu điện áp hiệu dụng hai đầu cuộn thứ cấp để hở là 40 V thì số vòng của cuộn này là
	\choice
	{300 vòng}
	{\True 200 vòng}
	{250 vòng}
	{400 vòng}
	\loigiai{
%\Binhluan{
		Xét cuộn thứ cấp: 
		\begin{align*}
		&E_{02}=U_2\sqrt{2}=40\sqrt{2}\ V.\\ 
		&E_{02}=\omega N_2\Phi_0\Rightarrow N=\dfrac{E_{02}}{\omega \Phi_0}=\dfrac{40\sqrt{2}}{100\pi. 0,9.10^{-3}}=200\ \ct{vòng}.
		\end{align*}
%}
%{
%Từ thông gửi qua một vòng dây: 
%\begin{align*}
%	\Phi_1=\Phi_2=\Phi \Leftrightarrow \Phi_{01}=\Phi_{02}=\Phi_0.
%\end{align*}
%Đề bài cho từ thông gửi qua tiết diện lõi sắt chính là từ thông gửi qua \textbf{một} vòng dây.
%}
	}
\end{vidu}
\loai{Viết biểu thức e}
\begin{vidu} %[3P3-14.2B][Loai6] 
	Từ thông xuyên qua một vòng dây của cuộn sơ cấp máy biến áp lí tưởng có dạng $\Phi =2\cos 100\pi t\ mWb$. Cuộn thứ cấp của máy biến áp có 1000 vòng. Biểu thức suất điện động ở cuộn thứ cấp là
	\choice
	{$e=200\pi \cos \left( 100\pi t \right)\ V$}
	{\True $e=200\pi \cos \left( 100\pi t-0,5\pi
		\right)\ V$}
	{$e=100\pi \cos \left( 100\pi t+0,5\pi
		\right)\ V$}
	{$e=100\pi \cos \left( 100\pi t \right)\ V$}
	\loigiai{
%\Binhluan{
		\begin{itemize}
			\item Ta có: $\Phi_1=\Phi_2=\Phi$.
			\item Biểu thức suất điện động ở cuộn thứ cấp là:
			\begin{align*}
				e_2=-N_2\Phi'&=1000. 200\pi {{. 10}^{-3}}\cos \left( 100\pi t-\dfrac{\pi}{2} \right)\ V\\
				&=200\pi \cos \left( 100\pi t-\dfrac{\pi}{2} \right)\ V.
			\end{align*}
		\end{itemize}	
%	}
%{$\Phi$ là từ thông qua \textbf{một} vòng dây}
}
\end{vidu}
\loai{Mạch có tải}
\begin{vidu} %[3P3-14.2B][Loai7] 
	Một máy tăng thế lí tưởng có tỉ số vòng dây giữa các cuộn sơ cấp $N_1$ và thứ cấp $N_2$ là 3. Biết cường độ dòng điện trong cuộn sơ cấp và điện áp hiệu dụng giữa hai đầu cuộn sơ cấp lần lượt là ${I_1}=6\ A$ và ${U_1}=120\ V$. Cường độ dòng điện hiệu dụng trong cuộn thứ cấp và điện áp hiệu dụng giữa hai đầu cuộn thứ cấp lần lượt là
	\choice
	{\True 2 A và 360 V}
	{18 V và 360 V}
	{2 A và 40 V}
	{18 A và 40 V}
	\loigiai{
%		\Binhluan{
		Máy tăng thế nên ${N_1}<{N_2}$
		\begin{align*}  &\dfrac{N_1}{N_2}=\dfrac{1}{3} ;\
			{U_1}=120\ V ;\
			{I_1}=6 \ A\\ 
			\Rightarrow\  &\dfrac{N_1}{N_2}=\dfrac{U_1}{U_2}=\dfrac{I_2}{I_1}
			\Rightarrow  {U_2}=\dfrac{N_2}{N_1}.{U_1}=360\ V ;\
			{I_2}=\dfrac{N_1}{N_2}.{I_1}=2 \ A.
		\end{align*}
%	}
%{
%Học sinh cần bám vào dữ kiên máy tăng thế để biện luận tỉ số $\dfrac{N_1}{N_2}=\dfrac{1}{3}$.
%}
	}
\end{vidu}